\chapter{Tempered Distributions}
\label{ch:iii18}

Tempered distributions are continuous linear functionals on Schwartz space. They include polynomial-growth functions, measures, Dirac masses, and their derivatives while remaining stable under the Fourier transform.

\section*{Learning goals}
After this chapter, the reader should be able to:
\begin{itemize}
\item verify that a distribution is tempered by establishing polynomial-growth bounds on Schwartz test functions;
\item embed polynomially growing functions and finite measures into tempered distributions when justified;
\item distinguish general distributions from tempered distributions by their admissible test-function spaces;
\item apply differentiation and polynomial multiplication inside the tempered-distribution framework;
\item construct examples of tempered singular distributions such as Dirac masses and their derivatives;
\item identify growth behavior that prevents a distribution from being tempered;
\end{itemize}
\section*{Conceptual roadmap}
\[
\boxed{\text{structure}\;\longrightarrow\;\text{estimate}\;\longrightarrow\;\text{limit or representation}\;\longrightarrow\;\text{application}.}
\]

\paragraph{Mathematical setting.}
Pairings on $\mathcal S$ are linear in the test function, without conjugation. Polynomial-growth functions are understood to be measurable and locally integrable, with an a.e. polynomial bound outside a compact set. A measure with $|\mu|(B_R)\le C(1+R)^N$ and locally finite total variation is tempered, by summing its action over dyadic shells against a Schwartz weight of degree greater than $N$.

\section{Tempered distributions}
A tempered distribution is a continuous linear functional on $\mathcal S(\mathbb R^n)$.

\section{Regular tempered distributions}
Functions of at most polynomial growth define tempered distributions when paired with rapidly decaying Schwartz functions.


% BEGIN VOL03-EXPANSION III18-example-01
\begin{example}[Polynomial growth is absorbed by Schwartz decay]\label{ex:iii18-expand-01}
Let \(f(x)=(1+|x|)^N\). For a Schwartz test function \(\phi\), choose
\(M>N+n+1\). Then
\[
|f(x)\phi(x)|\le C(1+|x|)^{N-M}p_M(\phi),
\]
where \(p_M(\phi)=\sup_x(1+|x|)^M|\phi(x)|\). The remaining weight is
integrable, so
\[
\left|\int f\phi\right|\le C_M p_M(\phi).
\]
Thus polynomial growth defines a continuous tempered distribution.
\end{example}
% END VOL03-EXPANSION III18-example-01
\section{Schwartz topology}
Continuity is measured against finitely many weighted derivative seminorms.

\section{Differentiation}
Distributional derivatives preserve temperedness.


% BEGIN VOL03-EXPANSION III18-example-02
\begin{example}[Differentiating a tempered distribution]\label{ex:iii18-expand-02}
For \(T\in\mathcal S'\),
\[
\partial^\alpha T(\phi)=(-1)^{|\alpha|}T(\partial^\alpha\phi).
\]
Since differentiation is a continuous endomorphism of Schwartz space, the map
\(\phi\mapsto T(\partial^\alpha\phi)\) is continuous. Therefore every
distributional derivative of a tempered distribution is again tempered.
\end{example}
% END VOL03-EXPANSION III18-example-02
\section{Polynomial multiplication}
Multiplication by coordinate polynomials preserves temperedness.

\section{Dirac family}
Dirac masses and all derivatives are tempered.

\section{Measures}
Many measures with polynomial growth define tempered distributions.

\section{Fourier readiness}
The Fourier invariance of Schwartz space lets the Fourier transform be transferred to its dual.


% BEGIN VOL03-EXPANSION III18-example-03
\begin{example}[Why Fourier transform extends to the dual]\label{ex:iii18-expand-03}
The Fourier transform is a continuous bijection
\(\mathcal S\to\mathcal S\). Thus for \(T\in\mathcal S'\), one may define
the transformed functional by
\[
\widehat T(\phi)=T(\widehat\phi)
\]
with no inverse transform or conjugation in this linear dual pairing. Continuity of
the Schwartz transform immediately makes \(\widehat T\) tempered. The test
space was chosen precisely to make this dual transfer possible.
\end{example}
% END VOL03-EXPANSION III18-example-03
\section{Core structural results}

\begin{theorem}[Polynomial-growth functions are tempered]\label{thm:iii18-01}
If $|f(x)|\le C(1+|x|)^N$ and $f$ is locally integrable, then $f$ defines a tempered distribution.
\begin{proof}
Choose a Schwartz seminorm with decay power larger than $N+n$; it dominates the integral of $|f\phi|$ and yields continuity.
\end{proof}
\end{theorem}

\begin{theorem}[Derivatives preserve temperedness]\label{thm:iii18-02}
If $T\in\mathcal S'$, then $\partial^\alpha T\in\mathcal S'$.
\begin{proof}
The map $\phi\mapsto\partial^\alpha\phi$ is continuous on Schwartz space, and composing it with the continuous functional $T$ preserves continuity.
\end{proof}
\end{theorem}

\begin{theorem}[Polynomial multiplication preserves temperedness]\label{thm:iii18-03}
If $T\in\mathcal S'$ and $p$ is a polynomial, then $pT\in\mathcal S'$.
\begin{proof}
Multiplication by a polynomial is a continuous endomorphism of Schwartz space, so $pT(\phi)=T(p\phi)$ is continuous.
\end{proof}
\end{theorem}

\section{Worked examples}

\begin{example}[Polynomial growth gives a tempered distribution]\label{ex:iii18-worked-01}
If \(f\) satisfies
\[
|f(x)|\le C(1+|x|)^N,
\]
then for \(\varphi\in\mathcal S\),
\[
\left|\int f\varphi\right|
\le C\int(1+|x|)^N|\varphi(x)|dx.
\]
Choose a Schwartz seminorm with decay exponent \(N+n+1\); the remaining power
is integrable. Hence \(f\) defines a continuous functional on \(\mathcal S\).
\end{example}

\begin{example}[Dirac mass is tempered]\label{ex:iii18-worked-02}
For \(\varphi\in\mathcal S\),
\[
|\delta_a(\varphi)|=|\varphi(a)|
\le \sup_x|\varphi(x)|.
\]
Thus every point mass, and similarly every derivative of a point mass, is a
tempered distribution.
\end{example}

\begin{example}[A rapidly growing function need not be tempered by integration]\label{ex:iii18-worked-03}
The function \(e^{x^2}\) grows faster than every polynomial. Pairing it
with an arbitrary Schwartz function need not be integrable; for instance a
Gaussian with insufficient decay produces divergence. Polynomial growth is a
natural sufficient condition for regular tempered distributions.
\end{example}

\begin{example}[Differentiation preserves temperedness]\label{ex:iii18-worked-04}
If \(T\in\mathcal S'\), define
\[
\langle D^\alpha T,\varphi\rangle=(-1)^{|\alpha|}
\langle T,D^\alpha\varphi\rangle.
\]
Differentiation is continuous on \(\mathcal S\), so the composition remains a
continuous linear functional. Hence \(D^\alpha T\in\mathcal S'\).
\end{example}

\section{Solved dossiers}
Each dossier combines several ideas from the chapter in a longer argument. The aim is to make hypotheses, calculations, and proof strategy explicit.

\begin{problem}[Polynomial growth gives a tempered distribution]\label{prob:iii18-dossier-01}
Prove a polynomial-growth function defines a tempered distribution.
\end{problem}
\begin{solution}
Use rapid decay of Schwartz functions to dominate the polynomial growth and
bound the pairing by one Schwartz seminorm.
\end{solution}

\begin{problem}[Dirac mass is tempered]\label{prob:iii18-dossier-02}
Show \(\delta_a\in\mathcal S'\).
\end{problem}
\begin{solution}
Evaluation is bounded by the zeroth Schwartz seminorm.
\end{solution}

\begin{problem}[A rapidly growing function need not be tempered by integration]\label{prob:iii18-dossier-03}
Why does super-polynomial growth obstruct defining a regular tempered distribution by direct integration?
\end{problem}
\begin{solution}
Schwartz decay is faster than every polynomial but not faster than every
exponential; an exponentially growing function can defeat the decay of some
Schwartz tests.
\end{solution}

\begin{problem}[Differentiation preserves temperedness]\label{prob:iii18-dossier-04}
Prove derivatives of tempered distributions are tempered.
\end{problem}
\begin{solution}
The derivative operator maps \(\mathcal S\) continuously to itself, and
precomposition of a continuous functional with a continuous map remains
continuous.
\end{solution}

\begin{problem}[Constant function]\label{prob:iii18-dossier-05}
Show the constant function \(1\) defines a tempered distribution.
\end{problem}
\begin{solution}
It has polynomial growth of degree zero, and
\(|\int\varphi|\) is controlled by a suitable weighted Schwartz seminorm.
\end{solution}

\begin{problem}[Polynomial functions]\label{prob:iii18-dossier-06}
Why does every polynomial define a tempered distribution?
\end{problem}
\begin{solution}
A polynomial has polynomial growth, so the regular-distribution estimate
applies.
\end{solution}

\begin{problem}[Embedding of compactly supported distributions]\label{prob:iii18-dossier-07}
Why is every compactly supported distribution tempered?
\end{problem}
\begin{solution}
On one compact support, finitely many sup norms of derivatives control its
action; these are bounded by corresponding Schwartz seminorms.
\end{solution}

\begin{problem}[Schwartz duality]\label{prob:iii18-dossier-08}
What is the test-function space for tempered distributions?
\end{problem}
\begin{solution}
The Schwartz space \(\mathcal S(\mathbb R^n)\); tempered distributions are
its continuous linear functionals.
\end{solution}

\begin{problem}[Polynomial-growth functions are tempered proof dossier]\label{prob:iii18-dossier-09}
Prove the following structural statement and identify the step where its main hypothesis is used: If $|f(x)|\le C(1+|x|)^N$ and $f$ is locally integrable, then $f$ defines a tempered distribution.
\end{problem}
\begin{solution}
Choose a Schwartz seminorm with decay power larger than $N+n$; it dominates the integral of $|f\phi|$ and yields continuity. This proof also explains why the stated hypothesis is structural rather than cosmetic.
\end{solution}

\begin{problem}[Derivatives preserve temperedness proof dossier]\label{prob:iii18-dossier-10}
Prove the following structural statement and identify the step where its main hypothesis is used: If $T\in\mathcal S'$, then $\partial^\alpha T\in\mathcal S'$.
\end{problem}
\begin{solution}
The map $\phi\mapsto\partial^\alpha\phi$ is continuous on Schwartz space, and composing it with the continuous functional $T$ preserves continuity. This proof also explains why the stated hypothesis is structural rather than cosmetic.
\end{solution}

\begin{problem}[Polynomial multiplication preserves temperedness proof dossier]\label{prob:iii18-dossier-11}
Prove the following structural statement and identify the step where its main hypothesis is used: If $T\in\mathcal S'$ and $p$ is a polynomial, then $pT\in\mathcal S'$.
\end{problem}
\begin{solution}
Multiplication by a polynomial is a continuous endomorphism of Schwartz space, so $pT(\phi)=T(p\phi)$ is continuous. This proof also explains why the stated hypothesis is structural rather than cosmetic.
\end{solution}

\begin{problem}[Chapter synthesis]\label{prob:iii18-dossier-12}
Connect the main definitions and structural theorems of this chapter into one proof strategy. Explain what object is controlled, what estimate or closure principle is used, and what conclusion becomes available.
\end{problem}
\begin{solution}
Tempered distributions are continuous linear functionals on Schwartz space. They include polynomial-growth functions, measures, Dirac masses, and their derivatives while remaining stable under the Fourier transform. The recurring strategy is to encode the object in the chapter's natural measurable, normed, Fourier, or distributional structure, establish the controlling estimate, and then pass to the desired limit or representation.
\end{solution}
\section{Exercises with complete solutions}
The shorter exercise layer remains separate from the solved dossiers.

\begin{exercise}\label{exr:iii18-01}
State and apply the central principle behind Tempered distributions. Give one immediate consequence in the setting of this chapter.
\end{exercise}
\begin{hint}
Temperedness means continuity on Schwartz space; estimate the pairing by a finite combination of Schwartz seminorms.
\end{hint}
\begin{solution}
A tempered distribution is a continuous linear functional on $\mathcal S(\mathbb R^n)$. As an immediate consequence, the same principle can be inserted into later convergence, approximation, transform, or weak-form arguments whenever its hypotheses are verified.
\end{solution}

\begin{exercise}\label{exr:iii18-02}
State and apply the central principle behind Regular tempered distributions. Give one immediate consequence in the setting of this chapter.
\end{exercise}
\begin{hint}
A polynomially growing function is tempered because rapid decay of test functions makes the pairing absolutely integrable.
\end{hint}
\begin{solution}
Functions of at most polynomial growth define tempered distributions when paired with rapidly decaying Schwartz functions. As an immediate consequence, the same principle can be inserted into later convergence, approximation, transform, or weak-form arguments whenever its hypotheses are verified.
\end{solution}

\begin{exercise}\label{exr:iii18-03}
State and apply the central principle behind Schwartz topology. Give one immediate consequence in the setting of this chapter.
\end{exercise}
\begin{hint}
Finite measures are tempered when their growth is controlled sufficiently to pair with Schwartz functions.
\end{hint}
\begin{solution}
Continuity is measured against finitely many weighted derivative seminorms. As an immediate consequence, the same principle can be inserted into later convergence, approximation, transform, or weak-form arguments whenever its hypotheses are verified.
\end{solution}

\begin{exercise}\label{exr:iii18-04}
State and apply the central principle behind Differentiation. Give one immediate consequence in the setting of this chapter.
\end{exercise}
\begin{hint}
Dirac masses are tempered because point evaluation is bounded by a Schwartz seminorm.
\end{hint}
\begin{solution}
Distributional derivatives preserve temperedness. As an immediate consequence, the same principle can be inserted into later convergence, approximation, transform, or weak-form arguments whenever its hypotheses are verified.
\end{solution}

\begin{exercise}\label{exr:iii18-05}
State and apply the central principle behind Polynomial multiplication. Give one immediate consequence in the setting of this chapter.
\end{exercise}
\begin{hint}
Differentiation preserves temperedness by moving derivatives to the Schwartz test function.
\end{hint}
\begin{solution}
Multiplication by coordinate polynomials preserves temperedness. As an immediate consequence, the same principle can be inserted into later convergence, approximation, transform, or weak-form arguments whenever its hypotheses are verified.
\end{solution}

\begin{exercise}\label{exr:iii18-06}
State and apply the central principle behind Dirac family. Give one immediate consequence in the setting of this chapter.
\end{exercise}
\begin{hint}
Multiplication by polynomials preserves temperedness because polynomial weights are built into the Schwartz seminorms.
\end{hint}
\begin{solution}
Dirac masses and all derivatives are tempered. As an immediate consequence, the same principle can be inserted into later convergence, approximation, transform, or weak-form arguments whenever its hypotheses are verified.
\end{solution}

\begin{exercise}\label{exr:iii18-07}
State and apply the central principle behind Measures. Give one immediate consequence in the setting of this chapter.
\end{exercise}
\begin{hint}
To show a distribution is not tempered, build Schwartz functions whose pairing grows faster than any finite seminorm bound permits.
\end{hint}
\begin{solution}
Many measures with polynomial growth define tempered distributions. As an immediate consequence, the same principle can be inserted into later convergence, approximation, transform, or weak-form arguments whenever its hypotheses are verified.
\end{solution}

\begin{exercise}\label{exr:iii18-08}
State and apply the central principle behind Fourier readiness. Give one immediate consequence in the setting of this chapter.
\end{exercise}
\begin{hint}
Keep the test-function space visible: a distribution may be valid on \(C_c^\infty\) yet fail to extend continuously to Schwartz space.
\end{hint}
\begin{solution}
The Fourier invariance of Schwartz space lets the Fourier transform be transferred to its dual. As an immediate consequence, the same principle can be inserted into later convergence, approximation, transform, or weak-form arguments whenever its hypotheses are verified.
\end{solution}


% BEGIN VOL03-EXPANSION III18-exercises-01
\section{Graded supplementary exercises}
These exercises extend the preserved foundational set and are graded by mathematical role.

\subsection*{Standard computations}

\begin{exercise}[Polynomial]\label{exr:iii18-expand-01}
Does x\textasciicircum{}5 define a tempered distribution?
\end{exercise}
\begin{hint}
It has polynomial growth.
\end{hint}
\begin{solution}
Yes.
\end{solution}

\begin{exercise}[Dirac]\label{exr:iii18-expand-02}
Is delta\_0 tempered?
\end{exercise}
\begin{hint}
Point evaluation is controlled by a Schwartz seminorm.
\end{hint}
\begin{solution}
Yes.
\end{solution}

\begin{exercise}[Derivative]\label{exr:iii18-expand-03}
If T is tempered, what about DT?
\end{exercise}
\begin{hint}
Differentiation preserves temperedness.
\end{hint}
\begin{solution}
DT is tempered.
\end{solution}

\begin{exercise}[Polynomial multiplier]\label{exr:iii18-expand-04}
If T is tempered, what about x\textasciicircum{}3 T?
\end{exercise}
\begin{hint}
Polynomial multiplication preserves the class.
\end{hint}
\begin{solution}
It is tempered.
\end{solution}

\begin{exercise}[Finite measure]\label{exr:iii18-expand-05}
Does every finite Borel measure define a tempered distribution?
\end{exercise}
\begin{hint}
Pair with bounded Schwartz functions.
\end{hint}
\begin{solution}
Yes; absolute action is bounded by total variation times sup norm.
\end{solution}

\subsection*{Proofs}

\begin{exercise}[Polynomial-growth proof]\label{exr:iii18-expand-06}
Prove a locally integrable polynomial-growth function is tempered.
\end{exercise}
\begin{hint}
Use a Schwartz seminorm with decay power larger than growth plus dimension.
\end{hint}
\begin{solution}
The weighted integrand becomes integrable and gives a seminorm bound.
\end{solution}

\begin{exercise}[Derivative stability]\label{exr:iii18-expand-07}
Prove derivatives preserve temperedness.
\end{exercise}
\begin{hint}
Compose T with continuous differentiation on Schwartz space.
\end{hint}
\begin{solution}
Continuity is preserved.
\end{solution}

\begin{exercise}[Polynomial multiplication]\label{exr:iii18-expand-08}
Prove pT is tempered for polynomial p.
\end{exercise}
\begin{hint}
Multiplication by p is continuous on Schwartz space.
\end{hint}
\begin{solution}
Compose that map with T.
\end{solution}

\begin{exercise}[Dirac continuity]\label{exr:iii18-expand-09}
Prove delta and its derivatives are tempered.
\end{exercise}
\begin{hint}
Point derivatives are controlled by Schwartz derivative seminorms.
\end{hint}
\begin{solution}
Each action is a continuous seminorm evaluation.
\end{solution}

\subsection*{Counterexamples and hypothesis testing}

\begin{exercise}[Exponential growth]\label{exr:iii18-expand-10}
Why does e\textasciicircum{}(x\textasciicircum{}2) not define a regular tempered distribution by integration over the real line?
\end{exercise}
\begin{hint}
Schwartz functions need not decay faster than every exponential.
\end{hint}
\begin{solution}
Polynomial-weight seminorms cannot control the pairing for all Schwartz tests.
\end{solution}

\begin{exercise}[General distribution not tempered]\label{exr:iii18-expand-11}
Why is every distribution not automatically tempered?
\end{exercise}
\begin{hint}
Tempered distributions act on the larger Schwartz test space, not just compact supports.
\end{hint}
\begin{solution}
Growth at infinity can destroy continuity on Schwartz space.
\end{solution}

\begin{exercise}[Locally integrable not enough]\label{exr:iii18-expand-12}
Why is local integrability alone insufficient for temperedness?
\end{exercise}
\begin{hint}
Schwartz test functions need not have compact support, so behavior at infinity matters.
\end{hint}
\begin{solution}
Behavior at infinity must also be controlled.
\end{solution}

\subsection*{Applications and investigations}

\begin{exercise}[Frequency-ready sources]\label{exr:iii18-expand-13}
Why are tempered distributions suited to whole-space PDE?
\end{exercise}
\begin{hint}
They include singular/polynomial-growth sources and admit Fourier transform.
\end{hint}
\begin{solution}
Constant-coefficient PDE can be transformed even for nonsmooth data.
\end{solution}

\begin{exercise}[Idealized signals]\label{exr:iii18-expand-14}
Why can impulses and polynomial trends coexist in S-prime?
\end{exercise}
\begin{hint}
Dirac masses and polynomial-growth functions are both tempered.
\end{hint}
\begin{solution}
The common framework handles both under Fourier operations.
\end{solution}

\subsection*{Challenge problems}

\begin{exercise}[Order estimate]\label{exr:iii18-expand-15}
Show delta\_0' is controlled by a first-derivative Schwartz seminorm.
\end{exercise}
\begin{hint}
Its action is -phi'(0).
\end{hint}
\begin{solution}
Absolute value is at most sup |phi'|.
\end{solution}

\begin{exercise}[Growth threshold]\label{exr:iii18-expand-16}
If |f(x)| <= C(1+|x|)\textasciicircum{}N, explain how large a decay seminorm is needed in dimension n.
\end{exercise}
\begin{hint}
Choose a weight exponent greater than N+n.
\end{hint}
\begin{solution}
Then (1+|x|)\textasciicircum{}(N-M) is integrable, giving continuity.
\end{solution}
% END VOL03-EXPANSION III18-exercises-01
\section*{Chapter summary}
The chapter now supplies a canonical theorem-and-dossier layer for subsequent measure, Fourier, distributional, Sobolev, and PDE arguments.
